\documentclass[12pt]{ai_conquer2026}
\addbibresource{references.bib}

\title{%
  \begin{center}
        \Large {AI VIET NAM -- AI Conquer 2026}\\[.5ex]  
        \Huge \textbf{Điền tên project (hướng Đi làm) vào đây}\\[.5ex]
  \end{center}
}

% Author --------------------------------------------------------------------
\posttitle{
\author{
\begin{minipage}{\textwidth}
\centering
{
\vspace{-1.5em}
Microwave / Tên học viên
}\
\end{minipage}
}
}
\date{}

% Header --------------------------------------------------------------------
\pagestyle{fancy}
\fancyhf{}
\lhead{\href{https://www.facebook.com/aivietnam.edu.vn}{\textcolor{aivnGreen}{\bfseries AI VIET NAM (AI Conquer 2026)}}}
\rhead{\href{https://aivietnam.edu.vn/}{\textcolor{aivnGreen}{\bfseries aivietnam.edu.vn}}}
\fancyfoot[C]{\thepage}
\renewcommand{\headrulewidth}{1.0pt}
\renewcommand{\headrule}{{\color{aivnGreen}\hrule width\headwidth height\headrulewidth \vskip-\headrulewidth}}

\begin{document}
\maketitle
\vspace{-5.5em}
\setlength{\epigraphwidth}{0.6\textwidth}

\section{Tổng quan báo cáo}

Phần này giới thiệu ngắn gọn toàn bộ project mà nhóm thực hiện. Người viết cần nêu rõ project giải quyết bài toán gì, vì sao bài toán đó cần đến AI, sản phẩm cuối cùng là gì, và hệ thống hoạt động tổng quát như thế nào.

Nội dung nên bao gồm:
\begin{itemize}
\item Tên project và mục tiêu chính của project.
\item Vấn đề thực tế mà project muốn giải quyết.
\item Lý do sử dụng AI, Machine Learning hoặc Deep Learning cho bài toán này.
\item Người dùng mục tiêu hoặc bối cảnh sử dụng sản phẩm.
\item Dữ liệu được sử dụng trong project.
\item Mô hình hoặc phương pháp chính được áp dụng.
\item Kết quả nổi bật của mô hình hoặc sản phẩm demo.
\end{itemize}

Đối với hướng đi làm, phần tổng quan cần cho thấy project không chỉ là một bài thử nghiệm mô hình, mà là một sản phẩm có thể chạy được, có dữ liệu đầu vào, có kết quả đầu ra, có đánh giá mô hình và có demo rõ ràng.

\newpage
\setcounter{tocdepth}{2}
\tableofcontents
\thispagestyle{fancy}

\newpage
\section{Định nghĩa bài toán}

\subsection{Bối cảnh vấn đề}

Hiện nay, các mô hình ngôn ngữ lớn (LLM) đã và đang được sử dụng trong nhiều tác vụ như: hỏi đáp, nghiên cứu và tìm kiếm tài liệu. Tuy nhiên, trong bối cảnh giáo dục, nơi tính chính xác, khả năng kiểm chứng và sự bám sát với tài liệu nguồn là yêu cầu tiên quyết, LLM vẫn tồn tại những giới hạn cố hữu. Nguyên nhân đến từ ba hạn chế: tri thức chỉ dừng lại ở thời điểm huấn luyện, khả năng sinh ra thông tin sai lệch hoặc
không có căn cứ (hiện tượng “ảo giác”), và đặc biệt là sự thiếu hụt thông tin về dữ liệu nội bộ của từng tổ chức, cá nhân.

Project này được xây dựng nhằm mô phỏng một hệ thống hỗ trợ học tập dựa trên do người dùng cung cấp. Thay vì để LLM trả lời hoàn toàn từ tri thức nội tại, hệ thống sử dụng kiến trúc retrieval-augmented generation (RAG) để truy xuất các đoạn nội dung liên quan từ tài liệu đã được nạp trước, sau đó đưa các ngữ cảnh này vào prompt để sinh phản hồi. Cách tiếp cận này giúp câu trả lời có cơ sở rõ ràng hơn, giảm rủi ro hallucination và hỗ trợ người học kiểm tra lại nguồn thông tin.

\subsection{Mục tiêu}

Mục tiêu của project là xây dựng một hệ thống học tập như một phiên bản NotebookLM đơn
giản, có khả năng tiếp nhận các tài liệu học tập và hỗ trợ người dùng tìm hiểu nội dung tài liệu
thông qua LLM. Cụ thể project hướng đến các mục tiêu:
\begin{itemize}
    \item \textbf{Hỏi đáp dựa trên tài liệu}: Hệ thống cần nhận câu hỏi và tài liệu từ người dùng, truy xuất nội dung liên quan, trả lời dựa trên ngữ cảnh được cung cấp và kèm theo thông tin trích dẫn nguồn.
    \item \textbf{Giao diện thân thiện}: Người dùng upload PDF và trò chuyện trực tiếp qua giao diện chat, câu trả lời hiển thị trích dẫn dạng \texttt{[1]}, \texttt{[2]}.
    \item \textbf{Đánh giá chất lượng RAG}: Xây dựng framework đánh giá riêng, đo cả chất lượng truy xuất (retrieval) lẫn chất lượng câu trả lời (generation) trên nhiều bộ dữ liệu hỏi đáp chuẩn, thay vì chỉ đánh giá cảm quan.
    \item \textbf{Kiến trúc đơn giản, dễ triển khai}: Backend và frontend tách biệt rõ ràng, đóng gói bằng Docker Compose, có thể chạy được cục bộ chỉ với vài lệnh, không yêu cầu hạ tầng phức tạp.
\end{itemize}


\subsection{Đầu vào của hệ thống}

Đầu vào của hệ thống gồm hai nhóm chính tương ứng với hai chức năng của bài toán:

\begin{itemize}
    \item \textbf{Đối với chức năng hỏi đáp:}
    \begin{itemize}
        \item Tập tài liệu PDF do người dùng tải lên trong session.
        \item Câu hỏi của người dùng dưới dạng văn bản ngôn ngữ tự nhiên.
    \end{itemize}
    
    \item \textbf{Đối với chức năng sinh flashcards:}
    \begin{itemize}
        \item Các nguồn nội dung được trích chọn từ tập tài liệu trong session.
        \item Thông tin ngữ cảnh liên quan đến nội dung tài liệu đã được xử lý trước đó.
    \end{itemize}
\end{itemize}

Các tài liệu PDF sau khi được tải lên sẽ được tiền xử lý, trích xuất văn bản, chia nhỏ thành các đoạn phù hợp cho truy xuất và sử dụng làm nguồn tri thức cho toàn bộ hệ thống.

\subsection{Đầu ra mong muốn}

Đầu ra của hệ thống cũng được chia thành hai loại chính:

\begin{itemize}
    \item \textbf{Câu trả lời cho câu hỏi của người dùng:} Đây là văn bản trả lời ngắn gọn, đúng trọng tâm, được sinh ra dựa trên nội dung có trong tài liệu PDF đã cung cấp. Mỗi câu trả lời cần đi kèm \textit{citation} hoặc thông tin dẫn nguồn để người dùng có thể kiểm chứng lại nội dung trong tài liệu gốc.
    
    \item \textbf{Bộ flashcards dạng câu hỏi - trả lời:} Đây là tập các thẻ học tập được tạo từ các nội dung quan trọng trong tài liệu. Mỗi flashcard cần thể hiện một đơn vị kiến thức rõ ràng, trong đó câu hỏi phải dễ hiểu và câu trả lời phải chính xác, ngắn gọn, bám sát tài liệu nguồn.
\end{itemize}

\subsection{Loại bài toán AI}

Xét theo bản chất kỹ thuật, project này thuộc nhóm bài toán \textbf{Retrieval-Augmented Question Answering (RAG-QA)} kết hợp với \textbf{text generation}.

Cụ thể:

\begin{itemize}
    \item Chức năng hỏi đáp trên tài liệu PDF là một bài toán \textbf{Retrieval-Augmented Question Answering}, trong đó hệ thống cần truy xuất các đoạn văn bản liên quan từ tài liệu trước khi sinh câu trả lời.
    
    \item Chức năng sinh flashcards là một bài toán \textbf{text generation}, trong đó mô hình sử dụng nội dung tài liệu làm ngữ cảnh để tạo ra các cặp câu hỏi - trả lời phục vụ học tập.
\end{itemize}

Do đó, đây có thể xem là một hệ thống AI xử lý ngôn ngữ tự nhiên theo hướng \textbf{chatbot hỏi đáp tài liệu kết hợp sinh nội dung học tập}.

\subsection{Tiêu chí thành công}

Hệ thống được xem là đạt yêu cầu khi thỏa mãn các tiêu chí sau:

\begin{itemize}
    \item \textbf{Trả lời bám tài liệu:} Câu trả lời phải phản ánh đúng thông tin có trong các tài liệu PDF đã tải lên, hạn chế tối đa việc suy diễn ngoài dữ liệu nguồn.
    
    \item \textbf{Có citation:} Câu trả lời cần đi kèm thông tin trích dẫn hoặc dẫn nguồn từ tài liệu để tăng độ tin cậy và khả năng kiểm chứng.
    
    \item \textbf{Sinh được flashcards có ý nghĩa:} Các flashcards được tạo ra phải rõ ràng, chính xác, không quá dài, và có giá trị hỗ trợ học tập thực tế.
    
    \item \textbf{Luồng demo chạy end-to-end:} Hệ thống cần chạy hoàn chỉnh toàn bộ pipeline, từ bước tải tài liệu PDF, tiền xử lý, truy xuất thông tin, sinh câu trả lời có citation, đến bước sinh flashcards từ nội dung tài liệu.
    
    \item \textbf{Kết quả có thể giải thích ở mức cơ bản:} Người dùng có thể hiểu được câu trả lời được tạo ra dựa trên phần nào của tài liệu, thay vì nhận một kết quả thiếu căn cứ.
\end{itemize}

Tóm lại, project được định nghĩa là một hệ thống AI hỗ trợ khai thác tri thức từ tài liệu PDF, trong đó đầu vào là tài liệu và câu hỏi của người dùng, đầu ra là câu trả lời có dẫn nguồn và bộ flashcards học tập. Đây là một bài toán kết hợp giữa truy xuất thông tin và sinh ngôn ngữ, với trọng tâm là tính bám sát tài liệu, khả năng kiểm chứng và tính ứng dụng trong demo thực tế.



\section{Dữ liệu sử dụng}

Project sử dụng hai nhóm dữ liệu chính: dữ liệu vận hành phục vụ trải nghiệm hỏi đáp và học tập trong từng phiên làm việc, và dữ liệu đánh giá phục vụ đo lường chất lượng hệ thống. Điểm quan trọng là hệ thống không dựa trên một kho tri thức cố định được xây sẵn từ trước, mà tạo không gian tri thức theo từng \textit{chat session} từ các tài liệu người dùng tải lên.

\subsection{Nguồn dữ liệu}

Các nguồn dữ liệu chính được tóm tắt trong Bảng~\ref{tab:data-sources}.

\begin{table}[H]
\centering
\small
\begin{tabular}{|p{2.8cm}|p{4.5cm}|p{1.8cm}|p{5.3cm}|}
\hline
\textbf{Nhóm dữ liệu} & \textbf{Nguồn} & \textbf{Định dạng} & \textbf{Vai trò} \\ \hline
Dữ liệu vận hành & Tài liệu do người dùng tải lên trong từng session & PDF & Làm nguồn tri thức để truy xuất ngữ cảnh, trả lời câu hỏi, tạo tóm tắt và sinh flashcards \\ \hline
Dữ liệu đánh giá & Các file trong \texttt{backend/evaluation/datasets} hoặc JSON do module đánh giá sinh ra & JSON & Lưu câu hỏi, đáp án chuẩn, ngữ cảnh chuẩn và metadata để benchmark retrieval, generation và citation \\ \hline
\end{tabular}
\caption{Các nguồn dữ liệu được sử dụng trong hệ thống}
\label{tab:data-sources}
\end{table}

Trong phiên bản hiện tại, dữ liệu vận hành tập trung vào \textbf{PDF có lớp văn bản}. Đây là lựa chọn phù hợp với bối cảnh học tập, vì phần lớn tài liệu đầu vào là giáo trình, ghi chú hoặc tài liệu tham khảo dưới dạng PDF.

\subsection{Các thực thể dữ liệu chính}

Thay vì chỉ xem dữ liệu là một file PDF, hệ thống tổ chức dữ liệu theo nhiều mức biểu diễn khác nhau. Các thực thể quan trọng được mô tả trong Bảng~\ref{tab:data-entities}.

\begin{table}[H]
\centering
\small
\begin{tabular}{|p{2.6cm}|p{2.5cm}|p{8.2cm}|}
\hline
\textbf{Thực thể} & \textbf{Định dạng} & \textbf{Nội dung chính} \\ \hline
Document & PDF + metadata & Tên tài liệu, đường dẫn lưu trữ, loại file, trạng thái xử lý, số trang và thời điểm tạo \\ \hline
ChatSession & Metadata & Phiên làm việc dùng để nhóm tài liệu, lịch sử hội thoại và phạm vi truy xuất \\ \hline
Chunk & Text + metadata & Đoạn văn bản tách từ PDF, gắn với \texttt{doc\_id}, số trang, nội dung chunk và \texttt{parent\_text} \\ \hline
Embedding & Vector float & Biểu diễn ngữ nghĩa 384 chiều của từng chunk để phục vụ semantic retrieval \\ \hline
Citation & Text + metadata & Thông tin phục vụ dẫn nguồn như tên tài liệu, số trang, đoạn snippet và chỉ số \texttt{[1]}, \texttt{[2]} \\ \hline
EvalSample & JSON object & Câu hỏi, đáp án chuẩn, danh sách ngữ cảnh chuẩn và metadata như \texttt{source\_pdf}, \texttt{page} \\ \hline
\end{tabular}
\caption{Các thực thể dữ liệu chính trong hệ thống}
\label{tab:data-entities}
\end{table}

Nhìn từ góc độ dữ liệu, một PDF ban đầu chỉ là tài liệu thô. Sau khi được xử lý, nó được biểu diễn thành các đơn vị nhỏ hơn như chunk, embedding và citation. Đây là nền tảng để hệ thống vừa truy xuất được ngữ cảnh liên quan, vừa giải thích được câu trả lời bằng dẫn nguồn cụ thể.

\subsection{Đặc điểm dữ liệu}

Dữ liệu của project có một số đặc điểm quan trọng sau:

\begin{itemize}
    \item \textbf{Tính động theo session:} mỗi phiên làm việc có tập tài liệu riêng, do đó không tồn tại một kho dữ liệu dùng chung cho mọi người dùng.
    \item \textbf{Quy mô không cố định:} số lượng tài liệu, số trang và số chunk thay đổi tùy vào hành vi upload của người dùng.
    \item \textbf{Phụ thuộc mạnh vào chất lượng PDF:} hệ thống hiện ưu tiên PDF có text layer; PDF bị mã hóa hoặc không có text layer sẽ gây lỗi hoặc không dùng được.
    \item \textbf{Cần metadata để truy vết:} số trang, tên tài liệu và mã tài liệu là các thông tin bắt buộc để hiển thị citation và giới hạn truy xuất trong đúng session.
    \item \textbf{Tách biệt dữ liệu vận hành và dữ liệu đánh giá:} dữ liệu người dùng tải lên phục vụ hỏi đáp thời gian thực, trong khi dữ liệu JSON trong module evaluation phục vụ benchmark có kiểm soát.
\end{itemize}

Vì vậy, phần dữ liệu của hệ thống không chỉ gồm file đầu vào, mà còn gồm toàn bộ các biểu diễn trung gian giúp nối từ tài liệu thô tới câu trả lời có dẫn nguồn.

\section{Tiền xử lý và chuẩn bị dữ liệu}

Mục tiêu của giai đoạn này là biến tài liệu PDF thô thành các đơn vị văn bản có thể biểu diễn bằng vector và sử dụng an toàn trong hệ thống RAG. Ở đây, phần trình bày tập trung vào góc nhìn \textit{preprocessing}, tức là dừng lại ở bước tạo ra chunk, embedding và metadata cần thiết; chưa bàn đến lưu trữ chỉ mục, API hay workflow hệ thống.

\subsection{Document Parsing}

Bước đầu tiên là tiếp nhận và đọc nội dung PDF. Ở tầng backend, hệ thống chỉ chấp nhận file có phần mở rộng \texttt{.pdf} và MIME type phù hợp. Sau khi file được lưu xuống thư mục upload, module parser sử dụng \textbf{PyMuPDF} (thông qua \texttt{fitz}) để mở tài liệu và trích xuất văn bản theo từng trang.

Một số ràng buộc được áp dụng ngay ở bước này:
\begin{itemize}
    \item PDF bị mã hóa \textbf{không} được hỗ trợ trong phiên bản hiện tại.
    \item Nếu tài liệu không có lớp văn bản và parser không đọc được nội dung, hệ thống xem đó là dữ liệu không khả dụng.
    \item Các trang không có nội dung văn bản sẽ không được đưa sang bước tiếp theo.
\end{itemize}

Điều này cho thấy parsing không chỉ là đọc file, mà còn là bước lọc sơ bộ để xác định tài liệu nào thực sự có thể chuyển thành tri thức dùng được cho RAG. Phiên bản hiện tại chưa tích hợp OCR; do đó chất lượng PDF đầu vào ảnh hưởng trực tiếp tới toàn bộ pipeline phía sau.

\subsection{Text Cleaning và chuẩn hóa văn bản}

Sau khi trích xuất được văn bản theo trang, hệ thống thực hiện làm sạch ở mức tối thiểu nhưng cần thiết để giữ tính ổn định cho các bước sau:

\begin{itemize}
    \item loại bỏ khoảng trắng đầu và cuối của văn bản trang;
    \item bỏ qua những trang rỗng sau khi làm sạch;
    \item chuẩn hóa khoảng trắng bằng cách gộp nhiều dấu cách, dấu xuống dòng và khoảng trắng liên tiếp thành một chuỗi văn bản gọn hơn.
\end{itemize}

Cách làm này đơn giản nhưng phù hợp với mục tiêu của project: giữ nguyên nội dung ngữ nghĩa chính, đồng thời giảm nhiễu do định dạng PDF gây ra. Việc làm sạch vừa đủ cũng giúp hạn chế việc chia chunk bị lệch vì nhiều ký tự xuống dòng hoặc khoảng trắng không cần thiết.

\subsection{Chunking}

Sau bước làm sạch, văn bản được chia thành các đoạn nhỏ hơn để phục vụ truy xuất. Trong cấu hình hiện tại của backend, hệ thống sử dụng \texttt{chunk\_size = 1600} và \texttt{chunk\_overlap = 250}. Việc chia chunk được thực hiện theo cơ chế \textit{sliding window} trên văn bản của từng trang.

Thiết kế này có ba lợi ích chính:
\begin{itemize}
    \item \textbf{Giữ ngữ cảnh cục bộ:} mỗi chunk đủ dài để chứa một đơn vị nội dung có nghĩa.
    \item \textbf{Giảm mất mát thông tin ở biên:} phần overlap giúp nội dung nằm ở ranh giới giữa hai chunk không bị cắt đứt hoàn toàn.
    \item \textbf{Giữ liên kết với nguồn gốc dữ liệu:} mỗi chunk vẫn mang theo số trang để hỗ trợ citation ở giai đoạn sau.
\end{itemize}

Kết quả của bước này là một danh sách các \texttt{TextChunk}, trong đó mỗi phần tử có tối thiểu ba thành phần: số trang, nội dung chunk và bản văn bản cha đi kèm.

\subsection{Embedding Preparation}

Khi đã có chunk văn bản, hệ thống chuẩn bị dữ liệu cho bước biểu diễn ngữ nghĩa bằng mô hình embedding \texttt{intfloat/e5-small-v2}. Cụ thể:

\begin{itemize}
    \item với dữ liệu tài liệu, mỗi chunk được thêm tiền tố \texttt{passage:} trước khi encode;
    \item vector đầu ra được chuẩn hóa về đơn vị độ dài để thuận tiện cho truy xuất theo cosine distance;
    \item hệ thống kiểm tra số chiều vector đầu ra và kỳ vọng kích thước cố định là 384 chiều.
\end{itemize}

Sau bước này, mỗi đơn vị dữ liệu đã có thể được xem như một phần tử tri thức có cấu trúc, ví dụ:

\begin{aivnoutputbox}[Đầu ra sau tiền xử lý]
\begin{verbatim}
{
  "source": "lecture-01.pdf",
  "page": 12,
  "chunk": "... nội dung văn bản đã làm sạch ...",
  "parent_text": "... văn bản gốc của chunk ...",
  "embedding": [0.021, -0.114, ...]
}
\end{verbatim}
\end{aivnoutputbox}

Đến đây, dữ liệu đã chuyển từ dạng PDF thô sang dạng văn bản phân đoạn có vector ngữ nghĩa và metadata đủ để dùng cho các bước phía sau.

\section{Phương pháp và mô hình sử dụng}

\subsection{Tổng quan phương pháp RAG}

Project sử dụng kiến trúc \textbf{Retrieval-Augmented Generation (RAG)} làm phương pháp trung tâm. Thay vì để mô hình ngôn ngữ trả lời hoàn toàn từ tri thức có sẵn trong tham số, hệ thống trước hết truy xuất các chunk liên quan từ tài liệu người dùng tải lên, sau đó mới đưa các chunk này vào prompt để sinh câu trả lời.

Cách tiếp cận này phù hợp với bài toán học tập trên tài liệu vì hai lý do. Thứ nhất, nguồn tri thức thay đổi theo từng session nên không thể giải quyết bằng một tập train cố định. Thứ hai, người dùng cần câu trả lời có thể kiểm chứng bằng tài liệu gốc, chứ không chỉ là một phản hồi nghe hợp lý.

\subsection{Embedding Model}

Mô hình embedding hiện tại là \texttt{intfloat/e5-small-v2}. Đây là mô hình pretrained gọn, phù hợp với bài toán semantic search trên văn bản. Trong code, hệ thống tách riêng hai kiểu biểu diễn:

\begin{itemize}
    \item \texttt{passage: ...} cho chunk tài liệu;
    \item \texttt{query: ...} cho câu hỏi người dùng.
\end{itemize}

Cách gắn tiền tố này bám theo cách sử dụng khuyến nghị của họ E5, giúp tăng tính nhất quán giữa vector truy vấn và vector tài liệu.

\subsection{Vector Database}

Trong phiên bản hiện tại, hệ thống sử dụng \textbf{PostgreSQL kết hợp pgvector} để lưu embedding và thực hiện truy xuất vector. Đây là lựa chọn phù hợp với phase hiện tại vì cùng một hệ quản trị có thể lưu cả dữ liệu quan hệ như \texttt{Document}, \texttt{ChatSession}, \texttt{ChatMessage} lẫn vector của từng chunk.

Việc sử dụng \texttt{pgvector} cũng giúp backend dễ kết hợp truy xuất vector với metadata của tài liệu trong cùng một truy vấn, đặc biệt là khi cần lấy ra tên tài liệu, số trang và nội dung snippet để phục vụ citation.

\subsection{Retrieval Strategy}

Ở bước truy xuất, câu hỏi của người dùng được embedding rồi so khớp với embedding của các chunk bằng \textbf{cosine distance}. Truy xuất được giới hạn trong các tài liệu thuộc đúng session hiện tại, tức là hệ thống không tìm kiếm trên toàn bộ mọi tài liệu đã từng tải lên.

Chiến lược này phù hợp với bài toán hỏi đáp trên tài liệu cá nhân vì nó giúp thu hẹp không gian tìm kiếm vào đúng tập tài liệu mà người dùng đang làm việc, từ đó giảm nhiễu và tăng tính liên quan của ngữ cảnh được chọn.

\subsection{Reranking}

Ngoài truy xuất dense ban đầu, backend còn hỗ trợ \textbf{reranking}. Cấu hình mặc định hiện tại dùng \texttt{BAAI/bge-reranker-base}, lấy trước một tập ứng viên rồi sắp xếp lại để chọn ra \textit{top-k} phù hợp hơn cho bước generation.

Thiết kế này giúp tăng chất lượng ngữ cảnh đầu vào mà không cần mở rộng quá nhiều số lượng chunk đưa vào prompt. Nói cách khác, retrieval đóng vai trò tìm nhanh các ứng viên gần đúng, còn reranking giúp chọn lại các chunk có giá trị nhất cho câu hỏi cụ thể.

\subsection{Prompt Strategy}

Prompt được xây dựng theo nguyên tắc \textbf{answer only from provided context}. Backend ghép các chunk đã truy xuất thành một khối ngữ cảnh có đánh số, bổ sung một phần \textit{recent conversation} ngắn để giữ mạch hội thoại, rồi chèn câu hỏi của người dùng ở cuối prompt.

Với chiến lược này, mô hình không được khuyến khích trả lời theo tri thức chung, mà phải ưu tiên nội dung trong các chunk đã cung cấp. Đây là một cơ chế quan trọng để giảm hallucination và tăng khả năng kiểm chứng câu trả lời.

\subsection{Large Language Model}

Tầng sinh câu trả lời hiện sử dụng \textbf{Gemini 2.5 Flash}. Đây là mô hình đảm nhận các tác vụ generation chính như trả lời câu hỏi, tạo tóm tắt và sinh nội dung học tập từ ngữ cảnh đã truy xuất.

Việc chọn Gemini 2.5 Flash phù hợp với mục tiêu của hệ thống vì mô hình có khả năng sinh ngôn ngữ tốt, tốc độ phản hồi phù hợp cho trải nghiệm chat và hỗ trợ cơ chế streaming để frontend hiển thị câu trả lời theo thời gian thực.

\subsection{Citation Mechanism}

Một yêu cầu quan trọng của hệ thống là mọi câu trả lời phải có khả năng truy vết về tài liệu nguồn. Vì vậy, mỗi chunk được đưa vào prompt đều được đánh số citation như \texttt{[1]}, \texttt{[2]}, và prompt cũng yêu cầu mô hình chèn citation inline ngay sau phát biểu có căn cứ.

Mỗi citation được gắn với metadata gồm:
\begin{itemize}
    \item chỉ số trích dẫn trong câu trả lời;
    \item mã tài liệu và tên tài liệu;
    \item số trang;
    \item snippet văn bản để frontend hiển thị khi hover.
\end{itemize}

Nhờ vậy, câu trả lời không chỉ đúng về nội dung mà còn có khả năng truy vết về đúng nguồn trong tài liệu.

Nhờ sự kết hợp giữa RAG, embedding, retrieval, reranking, prompt strategy và citation mechanism, hệ thống có thể tạo ra các câu trả lời bám tài liệu hơn so với việc chỉ gọi trực tiếp một mô hình ngôn ngữ lớn.


\section{Kiến trúc hệ thống}

\subsection{Kiến trúc tổng thể}

Kiến trúc hiện tại của project có thể mô tả ngắn gọn như sau:

\begin{aivnoutputbox}[Luồng kiến trúc mức cao]
Frontend (Next.js) 
$\downarrow$

Backend API (FastAPI)
$\downarrow$

Modules xử lý tài liệu + truy xuất + generation
$\downarrow$

PostgreSQL + pgvector \hspace{1em} Gemini 2.5 Flash
\end{aivnoutputbox}

Nếu nhìn riêng theo góc độ nạp tài liệu, luồng vận hành là:

\begin{aivnoutputbox}[Luồng nạp tài liệu]
Upload PDF
$\downarrow$

Parser
$\downarrow$

Chunker
$\downarrow$

Embedder
$\downarrow$

PostgreSQL + pgvector
\end{aivnoutputbox}

Trong phiên bản hiện tại, phần indexing chưa được tách thành một microservice độc lập; nó được thực hiện ngay trong backend thông qua module ingest. Đây là lựa chọn hợp lý cho phase đầu vì giúp hệ thống gọn, dễ chạy cục bộ và dễ debug.

\subsection{Vai trò của từng thành phần}

\begin{itemize}
    \item \textbf{Frontend:} giao diện Next.js chịu trách nhiệm tạo session, upload PDF, gửi câu hỏi, nhận answer streaming và hiển thị citation để người dùng kiểm tra nguồn.
    \item \textbf{Backend API:} FastAPI cung cấp các endpoint cho session, document, chat, summary và flashcards; đồng thời điều phối toàn bộ luồng xử lý tài liệu, truy xuất và sinh nội dung học tập.
    \item \textbf{Tầng dữ liệu:} PostgreSQL lưu tài liệu, chunk, lịch sử chat và metadata; phần \texttt{pgvector} lưu embedding để truy xuất theo cosine distance.
    \item \textbf{Tầng mô hình:} mô hình embedding và reranker chạy phía backend, còn Gemini 2.5 Flash đảm nhận generation.
    \item \textbf{Tầng chức năng học tập:} module study xây dựng summary, study notes và flashcards từ các nguồn nội dung đã truy xuất trong từng session.
    \item \textbf{Hạ tầng phụ trợ:} Docker Compose hiện khởi tạo PostgreSQL và Redis để phục vụ môi trường phát triển cục bộ.
\end{itemize}

\subsection{Flashcard generation trong kiến trúc hệ thống}

Chức năng flashcard phù hợp hơn khi được mô tả như một thành phần của kiến trúc hệ thống thay vì đặt trong chương mô hình. Trong luồng hiện tại, flashcard không phải là một mô hình độc lập được huấn luyện riêng, mà là một \textbf{khả năng của backend} được xây dựng trên cùng kho tri thức đã nạp vào session.

Cụ thể, backend lấy các citation hoặc nguồn học tập từ tài liệu đã truy xuất, gom nhóm theo document, tạo study notes trung gian rồi mới sinh ra các cặp \textbf{Question -- Answer}. Do đó, flashcard nên được xem là một lớp ứng dụng nằm phía trên retrieval và generation, tái sử dụng cùng embedding model, cùng cơ chế truy xuất và cùng LLM của hệ thống.
\subsection{Vì sao chọn PostgreSQL + pgvector và metadata}

Một quyết định kiến trúc quan trọng của project là dùng \textbf{PostgreSQL + pgvector} thay vì tách thêm một vector database riêng. Cách làm này có một số ưu điểm:

\begin{itemize}
    \item lưu được cả dữ liệu quan hệ (document, session, message) và vector embedding trong cùng một hệ;
    \item thuận tiện cho các phép join giữa chunk và tài liệu để dựng citation;
    \item giảm độ phức tạp khi triển khai vì chỉ cần quản lý một lớp lưu trữ chính cho dữ liệu ứng dụng.
\end{itemize}

Bên cạnh đó, \textbf{metadata} là thành phần bắt buộc chứ không phải phụ trợ. Nếu chỉ lưu embedding mà không lưu \texttt{doc\_id}, tên tài liệu, số trang hay snippet, hệ thống sẽ khó hiển thị citation, khó giới hạn truy xuất theo session và cũng khó giải thích câu trả lời cho người dùng.

\subsection{Redis, Docker và khả năng mở rộng}

Mặc dù cấu hình hiện tại đã có Redis trong file \texttt{docker-compose.yml} và trong phần settings, luồng ingest hiện tại vẫn được gọi trực tiếp trong API upload chứ chưa tách thành worker nền độc lập. Điều này cần được trình bày rõ trong báo cáo để tránh mô tả quá mức những gì hệ thống đang có.

Tuy nhiên, việc chuẩn bị sẵn Redis và Docker Compose vẫn có ý nghĩa:
\begin{itemize}
    \item \textbf{Docker Compose} giúp dựng nhanh môi trường cục bộ lặp lại được, đặc biệt là PostgreSQL và Redis.
    \item \textbf{Redis} là nền tảng phù hợp nếu mở rộng sang hàng đợi tác vụ hoặc caching trong các phiên bản tiếp theo.
    \item \textbf{Worker} là hướng mở rộng tự nhiên cho các tác vụ nặng như ingest PDF lớn, nhưng ở phiên bản hiện tại chưa phải thành phần hoàn thiện của hệ thống.
\end{itemize}

\section{Luồng xây dựng kho tri thức và luồng truy xuất/sinh câu trả lời}

Đối với một hệ thống RAG như project này, không có bước huấn luyện mô hình theo nghĩa truyền thống. Mô hình embedding, reranker và LLM đều là pretrained; hệ thống chủ yếu thực hiện \textbf{xây dựng kho tri thức} từ tài liệu người dùng tải lên và \textbf{truy xuất + sinh câu trả lời} tại thời điểm sử dụng.

\subsection{Luồng xây dựng kho tri thức}

Đây là luồng biến PDF thô thành kho tri thức có thể truy xuất được trong session:

\begin{enumerate}
    \item Người dùng tạo hoặc chọn một \textit{chat session}.
    \item Người dùng upload PDF vào endpoint gắn với session đó.
    \item Backend kiểm tra định dạng file, sinh tên lưu trữ an toàn và tạo bản ghi \texttt{Document}.
    \item Parser đọc PDF, tách văn bản theo từng trang và loại bỏ các trang không có text hữu ích.
    \item Văn bản được làm sạch và chia thành các chunk với kích thước và overlap đã cấu hình.
    \item Mỗi chunk được chuyển thành embedding vector bằng mô hình \texttt{intfloat/e5-small-v2}.
    \item Các chunk cùng embedding và metadata được ghi vào tầng lưu trữ; trạng thái tài liệu được cập nhật sang \texttt{ready} nếu xử lý thành công.
\end{enumerate}

Luồng này chính là quá trình \textbf{xây dựng chỉ mục tri thức} cho session hiện tại. Nó không làm thay đổi tham số mô hình, nên chính xác hơn khi gọi là \textit{knowledge base construction} thay vì \textit{training}.

\subsection{Luồng truy xuất và sinh câu trả lời}

Khi người dùng đặt câu hỏi, hệ thống vận hành theo trình tự sau:

\begin{enumerate}
    \item Frontend gửi câu hỏi của người dùng lên backend và lưu tin nhắn vào lịch sử hội thoại.
    \item Backend lấy một số tin nhắn gần nhất để giữ ngữ cảnh hội thoại ngắn hạn.
    \item Câu hỏi được embedding bằng cùng họ mô hình đã dùng cho tài liệu, nhưng với tiền tố \texttt{query:}.
    \item Hệ thống truy xuất \textit{top-k} chunk liên quan nhất trong phạm vi các tài liệu thuộc session hiện tại.
    \item Nếu bật reranker, danh sách ứng viên được sắp xếp lại trước khi đưa vào prompt.
    \item Backend dựng prompt gồm câu hỏi, recent conversation và các context chunk kèm chỉ số citation.
    \item Gemini 2.5 Flash sinh câu trả lời theo dạng streaming, với yêu cầu chèn citation inline như \texttt{[1]} hoặc \texttt{[1][2]}.
    \item Sau khi sinh xong, backend tính \textit{citation coverage}, lưu câu trả lời cùng danh sách citation và trả kết quả về frontend để hiển thị.
\end{enumerate}

Các chức năng học tập khác như summary hoặc flashcards tái sử dụng cùng kho tri thức đã xây dựng; điểm khác biệt chủ yếu nằm ở prompt và dạng đầu ra, chứ không thay đổi bản chất truy xuất ngữ cảnh từ tài liệu.

\section{Công cụ và công nghệ sử dụng}

Phần này liệt kê và giải thích các công cụ, thư viện, framework, nền tảng hoặc dịch vụ được sử dụng trong project AI. Người viết không nên chỉ liệt kê tên công nghệ, mà cần giải thích vai trò của từng công nghệ trong pipeline AI.

Nội dung nên bao gồm:

\begin{itemize}
\item \textbf{Ngôn ngữ lập trình:} Python là lựa chọn phổ biến cho AI project; có thể kết hợp JavaScript hoặc TypeScript nếu có frontend.
\item \textbf{Xử lý dữ liệu:} NumPy, Pandas, OpenCV, Pillow, librosa hoặc các thư viện xử lý dữ liệu khác.
\item \textbf{Machine Learning / Deep Learning:} Scikit-learn, PyTorch, TensorFlow, Keras hoặc Hugging Face Transformers.
\item \textbf{Model / Pretrained Model:} CNN, ResNet, EfficientNet, YOLO, BERT, CLIP, Sentence Transformers, LLM hoặc các mô hình khác.
\item \textbf{Vector Database nếu có:} FAISS, ChromaDB, Qdrant hoặc Pinecone cho các bài toán retrieval và RAG.
\item \textbf{Framework xây dựng demo:} Streamlit, Gradio, Flask, FastAPI, React hoặc Next.js.
\item \textbf{Quản lý mã nguồn:} Git và GitHub.
\item \textbf{Quản lý công việc:} Jira, Trello, GitHub Project hoặc AIO Board nếu có.
\item \textbf{Triển khai demo:} Hugging Face Spaces, Streamlit Cloud, Render, Vercel, Docker hoặc nền tảng cloud khác nếu có.
\end{itemize}

Với mỗi công nghệ quan trọng, \textbf{mới đối với module}, người viết nên nêu rõ để các team khác có thể nắm được:
\begin{itemize}
\item Công nghệ đó được dùng ở phần nào của project.
\item Vì sao nhóm chọn công nghệ đó.
\item Công nghệ đó giúp giải quyết vấn đề gì.
\item Có hạn chế hoặc lưu ý gì khi sử dụng hay không.
\end{itemize}




\begin{aivncodebox}[Ví dụ:]
Nhóm sử dụng PyTorch để huấn luyện mô hình vì framework này linh hoạt, phù hợp cho deep learning và dễ tùy chỉnh quá trình training. Streamlit được sử dụng để xây dựng demo vì có thể tạo giao diện nhanh và tích hợp trực tiếp với pipeline inference bằng Python.
\end{aivncodebox}

\section{Cấu trúc mã nguồn và hướng dẫn chạy}

Phần này mô tả cách tổ chức GitHub repository và hướng dẫn người đọc chạy lại project. Đây là phần thể hiện project có khả năng chạy lại và bàn giao cho người khác.

Người viết nên trình bày cấu trúc thư mục chính của project. Ví dụ:

\newpage
\begin{aivnoutputbox}[Cấu trúc thư mục tham khảo]
project-name/
|-- data/
|   |-- raw/
|   |-- processed/
|
|-- notebooks/
|   |-- exploratory_data_analysis.ipynb
|   |-- model_experiment.ipynb
|
|-- src/
|   |-- preprocessing.py
|   |-- train.py
|   |-- inference.py
|   |-- evaluate.py
|   |-- utils.py
|
|-- models/
|   |-- best_model.pt
|
|-- app/
|   |-- app.py
|
|-- reports/
|   |-- figures/
|
|-- requirements.txt
|-- README.md
|-- .gitignore

\end{aivnoutputbox}

Sau đó, giải thích ngắn gọn vai trò của từng thư mục:

\begin{itemize}
\item \texttt{data/}: chứa dữ liệu thô hoặc dữ liệu đã xử lý.
\item \texttt{notebooks/}: chứa notebook dùng để khám phá dữ liệu hoặc thử nghiệm mô hình.
\item \texttt{src/}: chứa source code chính cho preprocessing, training, inference và evaluation.
\item \texttt{models/}: chứa model checkpoint hoặc model đã huấn luyện.
\item \texttt{app/}: chứa code chạy demo AI.
\item \texttt{reports/}: chứa hình ảnh, biểu đồ, kết quả hoặc tài liệu báo cáo.
\item \texttt{requirements.txt}: liệt kê các thư viện cần cài đặt.
\item \texttt{README.md}: hướng dẫn tổng quan về project.
\end{itemize}

Phần hướng dẫn chạy project nên bao gồm các bước cơ bản:

\begin{aivncodebox}[Các câu lệnh hướng dẫn]
\begin{python}
git clone https://github.com/username/project-name.git
cd project-name
pip install -r requirements.txt

python src/preprocessing.py
python src/train.py
python src/evaluate.py
streamlit run app/app.py
\end{python}
\end{aivncodebox}

Nếu project không cần train lại model, người viết cần nói rõ người đọc chỉ cần tải checkpoint và chạy inference hoặc demo. Nếu project cần tải dữ liệu hoặc model checkpoint riêng, cần ghi rõ link tải và vị trí đặt file.

\section{Evaluation}

Phần này trình bày cách nhóm đánh giá chất lượng của hệ thống và sản phẩm demo. Đây là phần quan trọng để chứng minh hệ thống không chỉ chạy được, mà còn có kết quả đo lường cụ thể.

Đối với project này, các metric như retrieval metrics, generation metrics và citation-related metrics phù hợp hơn khi được đặt trong chương Evaluation thay vì chương phương pháp, vì đây là các công cụ dùng để đo hiệu quả của hệ thống sau khi đã thiết kế và triển khai.

\subsection{Các nhóm metric chính}

Module \texttt{backend/evaluation} đánh giá hệ thống ở cả hai mặt retrieval và generation. Các metric retrieval chính gồm:
\begin{itemize}
    \item \textbf{Recall};
    \item \textbf{Reciprocal Rank (RR / MRR ở mức mẫu)};
    \item \textbf{Precision@K};
    \item \textbf{Hit Rate@K};
    \item \textbf{NDCG@K};
    \item \textbf{Context Relevance} thông qua reranker.
\end{itemize}

Các metric generation chính gồm:
\begin{itemize}
    \item \textbf{Answer Similarity};
    \item \textbf{ROUGE-L};
    \item \textbf{Token F1};
    \item \textbf{Citation Coverage};
    \item \textbf{Faithfulness NLI}.
\end{itemize}

Nhờ bộ metric này, project không chỉ kiểm tra câu trả lời có trông hợp lý hay không, mà còn đo được mức độ truy xuất đúng ngữ cảnh, mức độ bám nguồn và độ nhất quán của câu trả lời với tài liệu đã cung cấp.

\subsection{Thiết lập dữ liệu đánh giá}

Dữ liệu đánh giá được tổ chức dưới dạng JSON trong module \texttt{backend/evaluation}, với mỗi mẫu gồm câu hỏi, đáp án chuẩn, ngữ cảnh chuẩn và metadata liên quan đến tài liệu nguồn. Thiết kế này giúp việc benchmark retrieval và generation được thực hiện có kiểm soát, tách biệt với dữ liệu vận hành mà người dùng tải lên trong quá trình demo.

Phần này trình bày cách nhóm đánh giá chất lượng của mô hình và sản phẩm demo. Đây là phần quan trọng để chứng minh mô hình không chỉ chạy được, mà còn có kết quả đo lường cụ thể.

Người viết nên trình bày:

\begin{itemize}
\item Dữ liệu test được sử dụng để đánh giá.
\item Metric đánh giá phù hợp với loại bài toán.
\item Baseline model hoặc phương pháp so sánh.
\item Kết quả của các phiên bản mô hình khác nhau.
\item Nhận xét về kết quả cuối cùng.
\end{itemize}

Tùy vào loại bài toán AI, có thể sử dụng các metric sau:

\begin{itemize}
\item \textbf{Classification:} Accuracy, Precision, Recall, F1-score, Confusion Matrix.
\item \textbf{Regression:} MAE, MSE, RMSE, R-squared.
\item \textbf{Object Detection:} IoU, mAP, Precision, Recall.
\item \textbf{Segmentation:} IoU, Dice Score, Pixel Accuracy.
\item \textbf{Retrieval / Search:} Precision@K, Recall@K, MRR.
\item \textbf{Generation / Chatbot:} ROUGE, BLEU, human evaluation, factual correctness hoặc relevance score.
\item \textbf{Product Demo:} Inference time, success rate, error rate, usability checklist.
\end{itemize}

Bên cạnh kết quả định lượng, người viết nên có phần nhận xét định tính. Ví dụ, mô hình xử lý tốt loại dữ liệu nào, còn yếu ở tình huống nào, và kết quả hiện tại đã đủ tốt để demo hay chưa.



\begin{aivncodebox}[Ví dụ:]
Mô hình cuối cùng đạt accuracy 86\% trên tập test, cao hơn baseline 12\%. Tuy nhiên, confusion matrix cho thấy mô hình vẫn nhầm lẫn giữa hai class có đặc điểm hình ảnh gần giống nhau. Điều này cho thấy nhóm cần bổ sung thêm dữ liệu hoặc cải thiện bước data augmentation.
\end{aivncodebox}

\section{Phân tích lỗi và hạn chế}

Phần này giúp báo cáo có chiều sâu hơn. Một technical report AI chuyên nghiệp không chỉ nêu kết quả tốt, mà còn cần phân tích các lỗi, hạn chế và nguyên nhân khiến mô hình chưa hoạt động như mong muốn.

Người viết nên trình bày:

\begin{itemize}
\item Những trường hợp mô hình dự đoán đúng hoặc xử lý tốt.
\item Những trường hợp mô hình dự đoán sai hoặc chưa ổn định.
\item Các class hoặc loại input mà mô hình thường nhầm lẫn.
\item Nguyên nhân có thể dẫn đến lỗi.
\item Hạn chế về dữ liệu, mô hình, quá trình huấn luyện, inference hoặc demo.
\item Những cải tiến có thể thực hiện nếu có thêm thời gian.
\end{itemize}



\begin{aivncodebox}[Ví dụ:]
Mô hình hoạt động tốt với ảnh rõ nét và ánh sáng đầy đủ, nhưng kết quả giảm khi ảnh bị mờ, thiếu sáng hoặc có nhiều vật thể gây nhiễu. Nguyên nhân có thể đến từ việc dữ liệu huấn luyện chưa đủ đa dạng và pipeline tiền xử lý chưa xử lý tốt các trường hợp ảnh chất lượng thấp.
\end{aivncodebox}

\newpage

\begin{aivncodebox}[Ví dụ:]
Nếu project là chatbot hoặc RAG, phần phân tích lỗi có thể tập trung vào:
\begin{itemize}
\item Truy xuất sai tài liệu.
\item Câu trả lời thiếu thông tin.
\item Câu trả lời đúng một phần nhưng diễn đạt chưa rõ.
\item Model trả lời ngoài phạm vi dữ liệu.
\item Chunking hoặc embedding chưa tối ưu.
\end{itemize}
\end{aivncodebox}

Phần này cho thấy nhóm không chỉ chạy được mô hình, mà còn hiểu điểm mạnh, điểm yếu và giới hạn hiện tại của hệ thống .

\section{Kết luận và hướng phát triển}

Phần này tổng kết lại toàn bộ project. Người viết nên nhắc lại bài toán ban đầu, mô hình hoặc phương pháp đã sử dụng, sản phẩm demo đã xây dựng, kết quả đạt được và những điểm có thể phát triển trong tương lai.

Nội dung kết luận nên trả lời các câu hỏi:

\begin{itemize}
\item Project đã giải quyết bài toán gì?
\item Dữ liệu và mô hình chính được sử dụng là gì?
\item Hệ thống hoạt động theo luồng nào?
\item Kết quả đánh giá mô hình đạt được là gì?
\item Demo cuối cùng có thể làm được những chức năng nào?
\item Nhóm đã học được gì sau quá trình thực hiện project?
\end{itemize}

Phần hướng phát triển nên nêu các cải tiến cụ thể, ví dụ:

\begin{itemize}
\item Thu thập thêm dữ liệu để cải thiện chất lượng mô hình.
\item Cải thiện data preprocessing hoặc data augmentation.
\item Thử nghiệm mô hình mạnh hơn hoặc fine-tune tốt hơn.
\item Tối ưu inference speed để hệ thống phản hồi nhanh hơn.
\item Cải thiện giao diện demo để người dùng dễ sử dụng hơn.
\item Thêm cơ chế giải thích kết quả mô hình.

\end{itemize}

Kết luận nên nhấn mạnh rằng project không chỉ là một thử nghiệm mô hình, mà là một sản phẩm demo có thể chạy được, có dữ liệu, có pipeline, có đánh giá và có khả năng tiếp tục phát triển.

\newpage
\section{Tài liệu tham khảo}
\label{section:references}
\nocite{*}
\vspace{-3.2em}
\printbibliography

\newpage
\begin{appendices}
\begin{enumerate}
\item \textbf{AIVIETNAM:} Thông tin về AIVIETNAM có thể đọc thêm tại \href{https://aivietnam.edu.vn/}{đây}.
\item \textbf{GitHub Repository:} Nhóm điền link GitHub repository của project tại đây.
\item \textbf{Demo Video:} Nhóm điền link video demo tại đây.
\item \textbf{Demo App / Website:} Nhóm điền link demo app hoặc website tại đây nếu có.
\item \textbf{Dataset:} Nhóm điền link dataset hoặc mô tả nguồn dữ liệu tại đây nếu có.
\item \textbf{Model Checkpoint:} Nhóm điền link model checkpoint tại đây nếu có.
\end{enumerate}
\end{appendices}

\end{document}
